%! Composition of Functions — Unit 6, Lesson 6.6 — Group Activity
\documentclass[10pt]{article}
\usepackage{algebra2-article}
\usepackage{algebra2-boxes}
\newcommand{\TallMath}[1]{$\displaystyle #1\rule[-1.4em]{0pt}{3.2em}$}
\newcommand{\lgrid}[4]{%
  \draw[step=1, linegray!40, very thin] (#1,#3) grid (#2,#4);
  \draw[->, semithick, charcoal] (#1,0)--(#2,0) node[below right, font=\scriptsize]{$x$};
  \draw[->, semithick, charcoal] (0,#3)--(0,#4) node[left, font=\scriptsize]{$y$};}

\begin{document}

\pageheader{Unit 6, Lesson 6.6}{Group Activity}
\namepartnerperiod

\vspace{0.08in}

\begin{headlinebox}{forestbg}
\small\textbf{Wired in Series.} Two machines, one after the other --- and the wire between them
carries exactly one number. With your partner, write that middle number down \emph{every single
time}. Most mistakes today are not arithmetic; they are running the machines in the wrong order.
\end{headlinebox}

\vspace{0.06in}

\begin{tcolorbox}[colback=white, colframe=black!40, title=\textbf{Tier R --- Remediate},
    fonttitle=\bfseries, arc=1.5mm, left=3mm, right=3mm, top=2mm, bottom=2mm, breakable]
\textbf{Part 1 --- One step at a time.} \; Let $f(x)=x+4$ and $g(x)=3x$.
\begin{center}
\renewcommand{\arraystretch}{1.8}
\begin{tabularx}{0.96\linewidth}{@{}l X X X@{}}
\hline
 & \textbf{Which runs first?} & \textbf{The middle number} & \textbf{Final answer} \\
\hline
$(f\circ g)(2)$ & \blank{1.6cm} & $g(2)=\blank{1.4cm}$ & $f(\blank{1.0cm})=\blank{1.4cm}$ \\
$(g\circ f)(2)$ & \blank{1.6cm} & $f(2)=\blank{1.4cm}$ & $g(\blank{1.0cm})=\blank{1.4cm}$ \\
\hline
\end{tabularx}
\end{center}
\vspace{-4pt}
Same two functions, same input, two answers. Which one is bigger, and why does the order do that?
\blank{6.0cm}

\vspace{6pt}
\textbf{Part 2 --- Now you run them.} Let $f(x)=2x-1$ and $g(x)=x^{2}$. Show the middle number.
\begin{center}
\renewcommand{\arraystretch}{1.8}
\begin{tabularx}{\linewidth}{@{}X X@{}}
(a) \; $(f\circ g)(3)$: \; middle $=\blank{1.2cm}$, \; answer $=\blank{1.4cm}$ &
(b) \; $(g\circ f)(3)$: \; middle $=\blank{1.2cm}$, \; answer $=\blank{1.4cm}$ \\
(c) \; $(f\circ g)(-2)$: \; middle $=\blank{1.2cm}$, \; answer $=\blank{1.4cm}$ &
(d) \; $(g\circ f)(0)$: \; middle $=\blank{1.2cm}$, \; answer $=\blank{1.4cm}$ \\
\end{tabularx}
\end{center}

\vspace{4pt}
\textbf{Part 3 --- Read it off the graphs.} These are the same two graphs from your notes.
\begin{center}
\begin{minipage}[c]{0.47\linewidth}
\centering
\begin{tikzpicture}[scale=0.34]
  \lgrid{-5}{7}{-3}{5}
  \draw[very thick, forest] (-4,-2)--(0,2)--(4,0)--(6,4);
  \fill[forest] (0,2) circle (3.4pt);
  \fill[forest] (4,0) circle (3.4pt);
  \node[forest, font=\scriptsize] at (-2.6,2.6) {$y=f(x)$};
\end{tikzpicture}
\end{minipage}
\hfill
\begin{minipage}[c]{0.47\linewidth}
\centering
\begin{tikzpicture}[scale=0.34]
  \lgrid{-5}{7}{-3}{5}
  \draw[very thick, navy] (-4,4)--(2,-2)--(6,2);
  \fill[navy] (2,-2) circle (3.4pt);
  \node[navy, font=\scriptsize] at (4.4,3.4) {$y=g(x)$};
\end{tikzpicture}
\end{minipage}
\end{center}
\vspace{-2pt}
\begin{center}
\renewcommand{\arraystretch}{1.7}
\begin{tabularx}{0.96\linewidth}{@{}X X X X@{}}
\hline
 & \textbf{Middle number} & \textbf{Read on which graph?} & \textbf{Answer} \\
\hline
(a) $(f\circ g)(-4)$ & \blank{1.4cm} & \blank{1.6cm} & \blank{1.4cm} \\
(b) $(g\circ f)(-4)$ & \blank{1.4cm} & \blank{1.6cm} & \blank{1.4cm} \\
(c) $(f\circ g)(6)$ & \blank{1.4cm} & \blank{1.6cm} & \blank{1.4cm} \\
(d) $(g\circ f)(2)$ & \blank{1.4cm} & \blank{1.6cm} & \blank{1.4cm} \\
\hline
\end{tabularx}
\end{center}
\vspace{-4pt}
Rows (a) and (b) start from the same input, $x=-4$. Explain to your partner why they do not finish
at the same place. \blank{6.0cm}
\end{tcolorbox}

\begin{tcolorbox}[colback=white, colframe=black!40, title=\textbf{Tier A --- Approaching Proficiency},
    fonttitle=\bfseries, arc=1.5mm, left=3mm, right=3mm, top=2mm, bottom=2mm, breakable]
\textbf{Part 1 --- Both orders, in general.} For each pair, find both compositions and then decide
whether they are the same function.
\begin{center}
\renewcommand{\arraystretch}{1.8}
\begin{tabularx}{\linewidth}{@{}l X X c@{}}
\hline
\textbf{$f$ and $g$} & \textbf{$(f\circ g)(x)$} & \textbf{$(g\circ f)(x)$} & \textbf{Same?} \\
\hline
(a) $f(x)=3x$, \; $g(x)=x-2$ & \blank{2.6cm} & \blank{2.6cm} & \blank{1.2cm} \\
(b) $f(x)=x^{2}$, \; $g(x)=x+1$ & \blank{2.6cm} & \blank{2.6cm} & \blank{1.2cm} \\
(c) $f(x)=\sqrt{x}$, \; $g(x)=x-9$ & \blank{2.6cm} & \blank{2.6cm} & \blank{1.2cm} \\
(d) $f(x)=x+3$, \; $g(x)=x-3$ & \blank{2.6cm} & \blank{2.6cm} & \blank{1.2cm} \\
(e) $f(x)=3x-1$, \; $g(x)=\dfrac{x+1}{3}$ & \blank{2.6cm} & \blank{2.6cm} & \blank{1.2cm} \\
\hline
\end{tabularx}
\end{center}
\vspace{1.35in}

\textbf{Part 2 --- What did you just find?}
\begin{enumerate}[label=\textbf{J\arabic*.}, leftmargin=2.4em, itemsep=6pt]
  \item Rows (d) and (e) behaved differently from the others: \emph{both} orders came out to the
        same thing, and that thing was remarkably plain. What was it? What are those two functions
        doing to each other --- and what word would you invent for a pair like that?
        \writelines{3}
  \item Row (c) is the one to be careful with. The two answers do not have the same domain. State
        both domains, and for each one say \emph{which stage} of the composition does the refusing.
        \writelines{3}
  \item A classmate answered row (b) with $(f\circ g)(x)=x^{2}+1$. That is a real function --- just
        not the one that was asked for. Which composition did they actually compute, and what is the
        one-sentence fix? \writelines{2}
\end{enumerate}

\vspace{4pt}
\textbf{Part 3 --- Domains, on purpose.} Find $(f\circ g)(x)$, simplify it, and \emph{then} work out
the domain from the two machines.
\begin{center}
\renewcommand{\arraystretch}{1.8}
\begin{tabularx}{\linewidth}{@{}l X X@{}}
\hline
\textbf{$f$ and $g$} & \textbf{$(f\circ g)(x)$, simplified} & \textbf{Domain} \\
\hline
(a) $f(x)=\sqrt{x}$, \; $g(x)=x+2$ & \blank{3.0cm} & \blank{3.0cm} \\
(b) $f(x)=\sqrt{x-1}$, \; $g(x)=4x$ & \blank{3.0cm} & \blank{3.0cm} \\
(c) $f(x)=x^{2}+1$, \; $g(x)=\sqrt{x}$ & \blank{3.0cm} & \blank{3.0cm} \\
\hline
\end{tabularx}
\end{center}
\vspace{-2pt}
Row (c) simplifies to something with no radical left in it at all. Why is its domain still
restricted? \writelines{2}
\end{tcolorbox}

\begin{tcolorbox}[colback=white, colframe=black!40, title=\textbf{Tier E --- Extension},
    fonttitle=\bfseries, arc=1.5mm, left=3mm, right=3mm, top=2mm, bottom=2mm, breakable]
\textbf{Part 1 --- Settle the jacket argument for every price.} Keep $f(x)=0.80x$ (the $20\%$
discount) and $g(x)=x-10$ (the coupon), but stop using \$$80$.
\begin{enumerate}[label=(\alph*), leftmargin=1.9em, itemsep=6pt]
  \item Find both compositions as rules. \; $(f\circ g)(x)=\blank{3.0cm}$ \; and \;
        $(g\circ f)(x)=\blank{3.0cm}$
        \vspace{0.35in}
  \item Subtract one from the other. The gap between the two orders is \blank{2.0cm} --- and it does
        not depend on the price at all. Explain where that number comes from in terms of the
        \$$10$ coupon and the $20\%$ discount. \writelines{2}
  \item Which order does the \emph{store} prefer, and which does the \emph{customer} prefer? Does
        your answer change for an expensive jacket? \writelines{2}
\end{enumerate}

\vspace{4pt}
\textbf{Part 2 --- When \emph{do} two linear functions commute?} Let $f(x)=ax+b$ and $g(x)=cx+d$.
\begin{enumerate}[label=(\alph*), leftmargin=1.9em, itemsep=6pt, start=4]
  \item Compose both ways and simplify. \; $(f\circ g)(x)=\blank{3.4cm}$, \;
        $(g\circ f)(x)=\blank{3.4cm}$
        \vspace{0.40in}
  \item The $x$-terms are identical no matter what. So the two compositions are equal for
        \emph{every} $x$ exactly when the constants match: \blank{4.0cm}
  \item Test your condition on Tier A row (d) ($f(x)=x+3$, $g(x)=x-3$) and on the jacket pair
        ($f(x)=0.8x$, $g(x)=x-10$). Does it predict the right verdict both times? \writelines{2}
\end{enumerate}

\vspace{4pt}
\textbf{Part 3 --- More than two machines.}
\begin{enumerate}[label=(\alph*), leftmargin=1.9em, itemsep=6pt, start=7]
  \item Let $f(x)=x^{2}$, $g(x)=x-1$, $h(x)=3x$. Find $(f\circ g\circ h)(2)$ by going right to left.
        \; middles: \blank{1.2cm}, \blank{1.2cm} \; answer $=\blank{1.4cm}$
  \item Take $h(x)=2\sqrt{x-5}+1$ apart into \emph{three} machines run in order. \\
        First: \blank{2.6cm} \quad Second: \blank{2.6cm} \quad Third: \blank{3.0cm}
  \item In Lesson 6.4 you called $2\sqrt{x-5}+1$ ``the parent, shifted right $5$, stretched by $2$,
        shifted up $1$.'' Compare that list with your answer to (h). What is a transformation,
        really? \writelines{2}
\end{enumerate}

\vspace{4pt}
\textbf{Part 4 --- The pair that almost undoes itself.} Complete both columns. Write ``undefined''
where the expression does not exist over the real numbers.
\begin{center}
\renewcommand{\arraystretch}{1.6}
\begin{tabularx}{0.9\linewidth}{@{}c X X@{}}
\hline
\textbf{$x$} & \TallMath{\left(\sqrt{x}\,\right)^{2}} & \TallMath{\sqrt{x^{2}}} \\
\hline
$-4$ & \blank{2.2cm} & \blank{2.2cm} \\
$-1$ & \blank{2.2cm} & \blank{2.2cm} \\
$0$ & \blank{2.2cm} & \blank{2.2cm} \\
$4$ & \blank{2.2cm} & \blank{2.2cm} \\
$9$ & \blank{2.2cm} & \blank{2.2cm} \\
\hline
\end{tabularx}
\end{center}
\begin{enumerate}[label=(\alph*), leftmargin=1.9em, itemsep=5pt, start=10]
  \item As formulas, $\left(\sqrt{x}\,\right)^{2}=\blank{1.2cm}$ and $\sqrt{x^{2}}=\blank{1.2cm}$
        (Lesson 6.1). One of them is defined for every real number; the other is not.
        \blank{5.0cm}
  \item On which inputs do \emph{both} columns hand you back the $x$ you started with?
        \blank{3.0cm}
  \item So squaring and square-rooting undo each other --- but only on part of the number line.
        Tomorrow you will meet the word for two functions that undo each other everywhere. What do
        you predict has to be done to $y=x^{2}$ first, to make that possible? \writelines{3}
\end{enumerate}
\end{tcolorbox}

\end{document}
